% !TeX root = ../papers/nc-p3.tex
\section{安定な非正則境界を通る全成分と非対称正則窓}
\label{sec:binary-component-crossing}
\begin{theorem}\label{thm:binary-component-crossing}
定理\ref{thm:binary-secant-obstructions}の有理標本および向きを固定した二元障害の一般点で、
二隣接切断階数6の全軌跡の軌道横断完備局所環は
\[
                  k[[x_1,\ldots,x_{13},h,a,b]]/(ah,bh)
\]
である（幾何点での剰余体を $k$ と書いた）。アフィンテンソルでは滑らかな軌道変数61個を加える。
二つの成分は滑らかで、添字保持基底変更を除いた次元は14と15、交叉は被約滑らかな13次元である。
前者は両対称部分で、後者の一般点はどちらの反対交換でも対称化できず、
有限被約点グラフを持つ最小周期4の標準quadratic Koszul左右AS正則窓である。
後者上の片側対称な二因子の一般点も同じ正則性と周期4を持つ。
一般の安定な二元障害点は、これらの正則な成分が交わるところで非正則に特殊化した点である。
\end{theorem}
\begin{proof}[四指標の法線と二次式]
標本の $\ell,c$ は互いに素で二行は独立なので、有限グラフでFrobeniusかつ非正則である。
定理\ref{thm:all-finite-windows-semistable}から安定で、そのアフィン軌道は61次元。
反対二交換の符号で周囲の線型空間を四指標 $++,+-,-+,--$ に分ける。
第一切断 $w=rGr^t$ で $N$ を $r^t$ の十次元核の基底列とする。
各指標の変分 $X$ に対し、両切断の一次条件は $N^tXN=0$ の一組と等しい。
二つのSchur行列を交換に関して和差に分ければ、各指標の式数は順に
$55,45,45,55$ である。実際一つの交換は二切断を交換し、両交換の積は行列を転置する。
同じ可逆な主六次小行列を両切断で選べばSchur補もこの作用と可換である。
標本での核次元は有理数上で正確に $45,16,16,0$ となる。
軌道接の各指標次元 $31,15,15,0$ を除くと、横断変数は自明指標14と両混合指標各1である。
軌道横断は線型補空間と基底変更の積で形式的に取れる。有限な二交換群に関し
軌道接の補空間も、消去する方程式と変数の組も同変に選べる。

陰関数定理で独立な一次式を解く。自明指標の式は全て消去される。
二つの混合指標の残る式は $af$ と $bg$ 各一本、最後の指標は $ab h_j$ の形である。
ここで $a,b$ が対応する混合変数で、$f,g,h_j$ は符号変換に不変。
$f,g$ の定数項は零で、両方とも自明指標方向での微分が非零となる。
この計算では一般逆行列 $D$ を同じ可逆主小行列の逆で取り、
自明指標接 $X$ と各混合接 $A$ に対して
\[
              N^t(XDA+ADX)N
\]
を対応する一次写像の一次元余核に射影する。両指標で非零であることを
\path{computations/verify_binary_component_crossing.py}が検算する。
全Jacobi階数を有理数と素数1009の両方で照合してから二次式を還元しているため、
剰余体での非零性は標数零の接・余核上での非零性を与える。
これは元のSchur形式の二次項であり、解いた一次変数からの補正は余核で消える。
従って $f,g$ はそれぞれ正則パラメータである。
\end{proof}
\begin{proof}[有限グラフの次元下界による高次式の決定]
近傍で基底変更を許した両対称部分は、横断上では $a=b=0$ である。
実際両対称なら反対位置の点座標を点ごとに移す射影線型写像が各側に存在する。
基点でその線型方程式の核は恒等写像の一次元だけである。
点座標の二次制限が単射なので、点上で比例する二つの線型列は全空間で比例するためである。
有限平坦な点代数上のこれらの方程式は近傍で書け、その核が存在すれば基点で選んだ小行列により
一意に正規化して可逆な同定を戻せる。両同定で基底を合わせれば文字通り交換不変である
（定理\ref{thm:finite-four-slot-rank}）。片側の交換についても同様で、
同定後のテンソルが比例するとき符号は基点の正符号のままである。
従って対称部分の飽和は局所的に固定部分に滑らかな対応軌道を加えたものとなる。
指標ごとに横断を取った上では、各片側の固定条件は対応する混合座標の消失である。

横断の16変数正則完備環において $(b,f)$ は高さ2の素イデアルで $a$ を含まない。
この素点は全方程式を満たし、両対称でない。ここを通る極小成分は
系\ref{cor:finite-component-dimension-bound}により横断次元15以上である。
この下界は形式環でも使える。同系の証明は有限平坦な辺のSchur恒等式を開集合上で消去するもので、
両対称部分の外で同じ高さの上界が働く。あるいは有限型体上の局所環の優秀性と
完備化の等次元性を使い、両対称成分の完備化が $a=b=0$ に含まれることを使ってもよい。
従って全方程式の極小素イデアルで $(b,f)$ に含まれるものは正則環で高さ1以下であり、
$af$ を含み $a$ を含まないため $(f)$ そのものである。
特に $bg$ と全ての $ab h_j$ が $f$ の倍数である。
$f$ は $a,b$ と互いに素で $g$ は正則パラメータなので、$g$ は $f$ と同伴、
全 $h_j$ は $f$ の倍数となる。これで全イデアルは $(af,bf)$ と正確に決まり、
$h=f$ と取れる。次元の異なる二成分と交叉の主張が従う。
\end{proof}
\begin{proof}[二元軌跡の一般点と各成分での全次数]
この交叉が二元障害を持つ一般点に延びることを示す。
任意の二元表示で $K=\ker B$ は、$P_2(x)\otimes P_3(y)$ の乗法核
$K_0=\Delta(x,y)(P_1(x)\otimes P_2(y))$ を含む。$\Delta$ は二変数組の対角式である。
従って $\bigwedge^2 K$ 内に
\[
 \Delta(x,y)\Delta(x',y')\Delta(x',y)\Delta(x,y')\Delta(y,y')
\]
という非零元がある。これは $x,x'$ に対称、$y,y'$ に交代するので、反対の混合指標の
$N^tXN$ の値を全て消す。従って常に少なくとも一方の混合余核が非零である。
上の局所模型の両対称部分では、これが起こるのは交叉 $h=0$ に限られ、そこでは両方とも非零となる。
同じ結論を近傍で代数的に読むには、Jacobi階数の小行列と
共通の二成分および交叉の滑らかさを用いればよい。両対称成分だけの滑らかな点では
追加の混合接は存在しないので、標本近傍の二元軌跡は交叉に含まれる。
階数および混合二次式の非退化性は、この交叉上の開条件である。
二元パラメータの既約性から幾何学的一般点で同じ局所模型を得る。

交叉は横断次元13である。有限両対称Frobenius窓の非正則部分は向きごとに次元12しかないので、
交叉の一般点は全次数で正則となる。従ってこの交叉を含む二成分、および15次元成分内で
$a=0$ または $b=0$ の各片側対称因子の幾何学的一般点も正則である。これは各内部次数の
相対bar複体を正則点から一般点へ移す既出の特殊化判定による。
被約有限グラフは二元パラメータの一般点ですでに成立するから全てのこれらの一般点にも成立する。

一般二元表示の十次形式の射影安定群は自明である（一般位置の十根の集合を保つ射影直線の変換は恒等）。
商の一意性から $Q$ 座標の作用も自明で、$P_2$ の作用をスカラーまで固定した後に
第四文字を $sz+p$ に変える作用は $(s-1)Y+T\mathbin{\lrcorner}p=0$ を満たす必要があり恒等である。
従って添字を保つ射影安定群は自明となる。この点で二移動の同定は標準の反転以外になく、
番号付き二十点には恒等として働く。近傍で二移動がある場合も同じ置換でなければならない。
実際有限エタールな番号付けの後、指定置換の座標変換は射影行列空間で固有に極限を取れる。
特異行列が二組の二十点を対応させることはない。核と像の射影空間に入る点数の和は
階数1または3なら $7+1$ 以下、階数2なら $3+3$ 以下だからである。
従って各極限は可逆で、非恒等置換は特殊点でも同じ置換を与えてしまう。
恒等置換による二移動は反対両配置を点ごとに射影的に同定し、両対称性を強制する。
両対称部分の外にある15次元成分およびその片側対称因子の一般点では二移動も一移動もなく、
四移動により最小周期4である。
\end{proof}
